\scp{\it{k}-initial words: **kusay / *kusor / {\#}kisor(o) / **kindoR / **kVndoR(a)}{14-03}
There is a chain of incrementally resemblant forms meaning ‘cuscus’
stretching from Sulawesi through Halmahera across to the Bird’s
Head of Papua (or possibly in the other direction).
These forms all begin with *k
and are mostly followed by high vowels.
Their distribution is shown in \frf{14-03}.
No reflexes are known to exist in any of our Wallacean subgroups,
except in less than a handful of languages in the \tsc{Seram-Tanimbar-Bomberai} subgroup.
Within this subgroup, reflexes of **kVndoR(a) are known in only
one branch, the \tsc{Eastern Islands} node.
Outside of EMP, there is only one (!) language known to support
the full form with the vowels reconstructed as **kandoRa (as
reconstructed by \citealp{Blust1982,Blust2012b}).
This one language is an anomaly on several counts (see \srf{14-03-04})
when compared to the forms in nearby languages.

\begin{figure}[h]\centering
	\includegraphics[width=\columnwidth]{kandoRa.pdf}
	\fcp{Distribution of **kusay / *kusor / {\#}kisor(o) / **kVndoR(a)}{14-03}
\end{figure}